\documentclass[11pt]{amsart}

\usepackage[T1]{fontenc}
\usepackage{lmodern}
\usepackage{microtype}
\usepackage{amsmath,amssymb,amsthm}
\usepackage{booktabs,array}
\usepackage[hidelinks]{hyperref}

\hypersetup{
  pdftitle={A z-Visibility Representation of K6,6 Minus a Perfect Matching},
  pdfsubject={Three-dimensional visibility representations of graphs},
  pdfkeywords={z-visibility representation, visibility graph,
    complete bipartite graph, perfect matching, counterexample}
}

\newtheorem{theorem}{Theorem}
\newtheorem{lemma}[theorem]{Lemma}

\title[A $z$-visibility representation of $K_{6,6}-M$]
{A \texorpdfstring{$z$}{z}-Visibility Representation of
\texorpdfstring{$K_{6,6}$}{K6,6} Minus a Perfect Matching}
\date{}
\subjclass[2020]{05C62, 68U05}
\keywords{$z$-visibility representation, visibility graph, complete bipartite
graph, perfect matching, counterexample, three-dimensional visibility}

\begin{document}

\begin{abstract}
We exhibit an integral-coordinate $z$-visibility representation of
$K_{6,6}$ minus a perfect matching.  This directly refutes
Conjecture~6.5 of Bose et al. (1994), restated in the 1998 journal
version of that report.  A complete finite certificate is included.
\end{abstract}

\maketitle

\section{Counterexample}

Bose et al.~\cite[Sec.~1.2]{BoseEtAl1998} represent vertices by pairwise
disjoint closed axis-parallel rectangles in horizontal planes.  Two rectangles
$A$ and $B$ are $z$-visible if there is a closed vertical cylinder of positive
length and positive radius whose two end disks are contained in $A$ and in $B$
and whose intersection with every other rectangle of the arrangement is empty;
the 1994 report words this as ``a closed cylinder $C$ of non-zero length and
radius such that the ends of $C$ are contained in $R_i$ and $R_j$, the axis of
$C$ is parallel to the $z$ axis, and the intersection of $C$ with any other
rectangle in the arrangement is empty'' \cite[Sec.~1]{BoseEtAl1994}.
Adjacency is required to hold if and only if the corresponding rectangles are
$z$-visible, and a graph admitting such a representation is called
\emph{VRR-representable} there.  Writing $K_{j,j}-M$ for $K_{j,j}$ minus a
perfect matching, the report proves that $K_{5,5}-M$ is VRR-representable
\cite[Theorem~6.3]{BoseEtAl1994} and then conjectures, in the same section,
that
\begin{quote}
``$K_{6,6}-M$ is not VRR-representable''
\end{quote}
\cite[Conjecture~6.5]{BoseEtAl1994}; the conjecture was restated in
\cite[Sec.~5]{BoseEtAl1998}.  It carries no further qualifier there, in
particular not the ``directed version of'' qualifier of the adjacent
\cite[Theorem~6.4]{BoseEtAl1994}.

Let
\[
  H=\{h_0,\ldots,h_5\},\qquad V=\{v_0,\ldots,v_5\},
\]
and let $G$ have edge set
\[
  E(G)=\{h_iv_j:0\le i,j\le5,\ i\ne j\}.
\]
Thus $G$ is $K_{6,6}$ minus the perfect matching
$\{h_iv_i:0\le i\le5\}$.

\begin{theorem}\label{thm:main}
The graph $G$ has a $z$-visibility representation.  Hence
Conjecture~6.5 of Bose et al.~
\cite[Conjecture~6.5]{BoseEtAl1994} is false.
\end{theorem}

For each rectangle $R$, the table gives its height $z(R)$ and its projection
$\pi(R)$ to the $xy$-plane; the rectangle itself is
$\pi(R)\times\{z(R)\}$.  The distinct heights make the rectangles pairwise
disjoint.

\begin{center}
\small
\begin{tabular}{@{}c r c@{}}
\toprule
$R$ & $z(R)$ & $\pi(R)$ \\
\midrule
$h_0$ & 0  & $[6,9]\times[0,7]$ \\
$h_3$ & 1  & $[2,5]\times[0,7]$ \\
$v_4$ & 2  & $[4,8]\times[1,3]$ \\
$v_0$ & 3  & $[2,4]\times[1,6]$ \\
$v_5$ & 4  & $[4,8]\times[3,6]$ \\
$h_1$ & 5  & $[1,8]\times[2,5]$ \\
$v_3$ & 6  & $[2,10]\times[2,5]$ \\
$h_2$ & 7  & $[3,7]\times[1,6]$ \\
$v_1$ & 8  & $[3,7]\times[1,7]$ \\
$h_5$ & 9  & $[2,8]\times[1,3]$ \\
$h_4$ & 10 & $[2,8]\times[3,6]$ \\
$v_2$ & 11 & $[0,8]\times[0,4]$ \\
\bottomrule
\end{tabular}
\end{center}

\section{Verification}

For integers $a,b$, set
\[
  Q_{a,b}=(a,a+1)\times(b,b+1),
  \qquad
  \mathcal C(R)=
  \{(a,b)\in\mathbb Z^2:Q_{a,b}\subseteq\pi(R)\}.
\]
For distinct rectangles $A,B$, let
\[
  \mathcal I(A,B)=
  \{R:\min\{z(A),z(B)\}<z(R)<\max\{z(A),z(B)\}\}
\]
and
\[
  \mathcal F(A,B)=
  \bigl(\mathcal C(A)\cap\mathcal C(B)\bigr)
  \setminus
  \bigcup_{R\in\mathcal I(A,B)}\mathcal C(R).
\]

\begin{lemma}[Unit-cell criterion]\label{lem:cells}
Let the rectangles of the arrangement lie at pairwise distinct heights, so that
they are pairwise disjoint, and let every projection boundary lie on an integer
grid line.  Then two of them, $A$ and $B$, are $z$-visible if and only if
$\mathcal F(A,B)\ne\varnothing$.
\end{lemma}

\begin{proof}
Let $D$ be the horizontal base disk of a witnessing cylinder.  Choose
$p\in\operatorname{int}D$ on no integer grid line, and let $Q_{a,b}$ be the
open unit cell containing $p$.  Every projection boundary lies on an integer
grid line, so for every rectangle $R$, either
$Q_{a,b}\subseteq\pi(R)$ or
$Q_{a,b}\cap\pi(R)=\varnothing$.  Since $D$ lies in both endpoint projections
and avoids every intermediate projection,
$(a,b)\in\mathcal F(A,B)$.

Conversely, if $(a,b)\in\mathcal F(A,B)$, the closed disk of radius $1/4$
centered at $(a+1/2,b+1/2)$ lies in $Q_{a,b}$.  Its vertical extrusion from
$A$ to $B$ is a witnessing cylinder.
\end{proof}

For each $i\ne j$, the following entry is an element of
$\mathcal F(h_i,v_j)$; a dash marks the deleted matching edge.

\begin{center}
\small
\renewcommand{\arraystretch}{1.15}
\begin{tabular}{@{}c|cccccc@{}}
 & $v_0$ & $v_1$ & $v_2$ & $v_3$ & $v_4$ & $v_5$ \\
\hline
$h_0$ & --      & $(6,6)$ & $(6,0)$ & $(8,2)$ & $(6,1)$ & $(6,3)$ \\
$h_1$ & $(2,2)$ & --      & $(1,2)$ & $(2,2)$ & $(4,2)$ & $(4,3)$ \\
$h_2$ & $(3,1)$ & $(3,1)$ & --      & $(3,2)$ & $(4,1)$ & $(4,5)$ \\
$h_3$ & $(2,1)$ & $(3,6)$ & $(2,0)$ & --      & $(4,1)$ & $(4,3)$ \\
$h_4$ & $(2,5)$ & $(3,3)$ & $(2,3)$ & $(2,3)$ & --      & $(7,5)$ \\
$h_5$ & $(2,1)$ & $(3,1)$ & $(2,1)$ & $(2,2)$ & $(7,1)$ & --      \\
\end{tabular}
\end{center}

For every remaining pair $A,B$, the next table lists rectangles in
$\mathcal I(A,B)$ whose cell sets collectively cover
$\mathcal C(A)\cap\mathcal C(B)$.  The symbol $\varnothing$ means that this
intersection of cell sets is empty.  Thus $\mathcal F(A,B)=\varnothing$ in
every listed case.

\begin{center}
\footnotesize
\renewcommand{\arraystretch}{1.12}
\begin{tabular}{@{}c l@{\qquad}c l@{\qquad}c l@{}}
\toprule
pair & blockers & pair & blockers & pair & blockers \\
\midrule
\multicolumn{6}{c}{Pairs within $H$} \\
$h_0,h_1$ & $v_4,v_5$
  & $h_0,h_2$ & $v_4,v_5$
  & $h_0,h_3$ & $\varnothing$ \\
$h_0,h_4$ & $v_5$
  & $h_0,h_5$ & $v_4$
  & $h_1,h_2$ & $v_3$ \\
$h_1,h_3$ & $v_0,v_4,v_5$
  & $h_1,h_4$ & $v_3$
  & $h_1,h_5$ & $v_3$ \\
$h_2,h_3$ & $v_0,v_4,v_5$
  & $h_2,h_4$ & $v_1$
  & $h_2,h_5$ & $v_1$ \\
$h_3,h_4$ & $h_2,v_0$
  & $h_3,h_5$ & $h_2,v_0$
  & $h_4,h_5$ & $\varnothing$ \\
\addlinespace
\multicolumn{6}{c}{Pairs within $V$} \\
$v_0,v_1$ & $h_2$
  & $v_0,v_2$ & $h_1,h_5$
  & $v_0,v_3$ & $h_1$ \\
$v_0,v_4$ & $\varnothing$
  & $v_0,v_5$ & $\varnothing$
  & $v_1,v_2$ & $h_4,h_5$ \\
$v_1,v_3$ & $h_2$
  & $v_1,v_4$ & $h_2$
  & $v_1,v_5$ & $h_2$ \\
$v_2,v_3$ & $h_4,h_5$
  & $v_2,v_4$ & $h_5$
  & $v_2,v_5$ & $h_1$ \\
$v_3,v_4$ & $h_1$
  & $v_3,v_5$ & $h_1$
  & $v_4,v_5$ & $\varnothing$ \\
\addlinespace
\multicolumn{6}{c}{Deleted matching pairs} \\
$h_0,v_0$ & $\varnothing$
  & $h_1,v_1$ & $h_2$
  & $h_2,v_2$ & $v_1$ \\
$h_3,v_3$ & $h_1$
  & $h_4,v_4$ & $\varnothing$
  & $h_5,v_5$ & $\varnothing$ \\
\bottomrule
\end{tabular}
\end{center}

Every entry of the two tables is verified by comparing integer coordinates.
For the first table: the listed cell lies in $\mathcal C(h_i)$ and in
$\mathcal C(v_j)$, and in $\mathcal C(R)$ for no $R\in\mathcal I(h_i,v_j)$.
For the second: each listed blocker has height strictly between the two heights
of its pair, and the listed blockers' cell sets together cover
$\mathcal C(A)\cap\mathcal C(B)$, a cover that needs more than one blocker in
nine of the rows; in the eight rows marked $\varnothing$ the set
$\mathcal C(A)\cap\mathcal C(B)$ is itself empty, so no blocker is needed.
By Lemma~\ref{lem:cells},
the first table certifies all $30$ edges of $G$, while the second excludes
the other $36$ pairs.  Since
\[
  30+36=\binom{12}{2},
\]
the visibility graph of the twelve rectangles is exactly $G$.  This proves
Theorem~\ref{thm:main}.

\medskip
\noindent
\textbf{Scope.}
What is settled is exactly \cite[Conjecture~6.5]{BoseEtAl1994}: it asserts that
no representation of $K_{6,6}-M$ exists, and one representation refutes it.
Nothing else is claimed.  We prove nothing about $K_{j,j}-M$ for $j\ge7$, and we
make no minimality claim: neither that twelve rectangles with these coordinate
ranges are in any sense smallest, nor that the certificate above is the
shortest possible.

The reading of the definition is load-bearing, and it is the source's own: the
end disks of the cylinder are \emph{contained} in the two rectangles.  Under
the weaker reading in which those disks need only meet the two rectangles,
these twelve rectangles realise a strictly larger graph and would not represent
$G$; so the containment reading is not a normalisation of convenience but the
definition being refuted.  The verification is, by contrast, insensitive to
whether ``meets another rectangle'' is read as closed or as interior
intersection, since each exhibited witness keeps clearance $1/4$ from every
intermediate projection and an empty $\mathcal F(A,B)$ admits no cylinder of
positive radius under either reading.

The heights above interleave the two sides, in the order
$h_0$, $h_3$, $v_4$, $v_0$, $v_5$, $h_1$, $v_3$, $h_2$, $v_1$, $h_5$, $h_4$,
$v_2$.  The representation is
therefore not a \emph{directed} one in the sense of
\cite[Sec.~6]{BoseEtAl1994}, where representing a directed graph additionally
requires the rectangle of each edge's source to lie below the rectangle of its
target, and nothing here bears on \cite[Theorem~6.4]{BoseEtAl1994}, that the
directed version of $K_{5,5}-M$ is not VRR-representable.

We have found no resolution of Conjecture~6.5 in the literature; it is restated
as an open problem in \cite[Sec.~5]{BoseEtAl1998}, the journal version of the
report, published four years later.  That is the outcome of a search, not a
proof of priority.  The 1994 report has no DOI and we found no library route to
it; the copy read for this note, and the source of the wording quoted in
Section~1, is the one at the address printed in \cite{BoseEtAl1994}.  A
neighbouring model in the same literature, box visibility, where the channel
may run in any of the three axis directions, is treated in the companion note
\cite{CompanionCubes}, which exhibits nine unit cubes whose box-visibility
graph is $K_9$ and so refutes Conjecture~26 of \cite{FeketeMeijer1999}.
Different source, different statement, different model: neither result bears on
the other.

\medskip
\noindent
\textbf{Reproducibility.}
The inputs are exactly the data printed above: the twelve rectangles, the two
certificate tables, the definitions of Section~2, and the combinatorial
description of $G$.  The counts $30$, $36$ and $66=\binom{12}{2}$, the
partition of the $66$ pairs between the two tables, every cell set, and the
visibility graph itself are recomputed rather than read off, and that the
result is $K_{6,6}$ minus a perfect matching is confirmed without the
labelling, from bipartiteness, $5$-regularity, and complementarity in
$K_{6,6}$ to a perfect matching.  The visibility graph is built three
ways---by the unit-cell criterion of Lemma~\ref{lem:cells}, by the exact area
of the free region over the arrangement of the printed coordinates, and by
inclusion--exclusion over $\mathcal I(A,B)$---and all three agree on the $66$
pairs; for each of the $30$ edges the closed disk of radius $1/4$ is checked in
exact rational arithmetic to lie in both endpoint projections and to miss every
intermediate one.  The check discriminates: over the full census of all $96$
admissible unit perturbations of a single projection edge and all $66$
transpositions of the height order, $79$ and $61$ respectively destroy the
representation, and the
three methods agree on all $162$ mutants.  All arithmetic is exact and nothing
was sampled or reduced in range.  The computation was re-run independently on a
separate machine, where all $60$ of its checks agreed; the program and its
transcript are retained in an auxiliary archive available on request.  Such
agreement is evidence about the computation, not a second proof of
Lemma~\ref{lem:cells}, whose conclusion it confirms on this instance by the two
methods that never use it.

\begin{thebibliography}{99}

\bibitem{BoseEtAl1994}
P.~Bose, H.~Everett, S.~P.~Fekete, A.~Lubiw, H.~Meijer, K.~Romanik,
T.~Shermer, and S.~Whitesides,
\emph{On a visibility representation for graphs in three dimensions},
in D.~Avis and P.~Bose (eds.),
\emph{Snapshots of Computational and Discrete Geometry}, vol.~3,
McGill University School of Computer Science Technical Report SOCS-94.50,
July 1994, pp.~2--25.  The report has no DOI; a copy is available at
\url{https://www.ibr.cs.tu-bs.de/users/fekete/hp/publications/PDF/1994-On_Visibility_Representation_for_Graphs_in_Three_Dimensions.pdf}.

\bibitem{BoseEtAl1998}
P.~Bose, H.~Everett, S.~P.~Fekete, M.~E.~Houle, A.~Lubiw, H.~Meijer,
K.~Romanik, G.~Rote, T.~C.~Shermer, S.~Whitesides, and C.~Zelle,
\emph{A visibility representation for graphs in three dimensions},
J. Graph Algorithms Appl. \textbf{2} (1998), no.~3, 1--16,
\href{https://doi.org/10.7155/jgaa.00006}{doi:10.7155/jgaa.00006}.

\bibitem{FeketeMeijer1999}
S.~P.~Fekete and H.~Meijer,
\emph{Rectangle and box visibility graphs in 3D},
Internat. J. Comput. Geom. Appl. \textbf{9} (1999), no.~1, 1--27,
\href{https://doi.org/10.1142/S0218195999000029}
{doi:10.1142/S0218195999000029}.

\bibitem{CompanionCubes}
\emph{A unit-cube box-visibility representation of $K_9$},
companion note in this collection, file
\texttt{nine-unit-cubes-\allowbreak disprove-fekete-meijer-conjecture}.

\end{thebibliography}

\end{document}